\section{Analytic excitation of conical points by one genuine resonance}
\label{hs:modal-next-section}

This section proves a local bridge that the driven-field construction
does not provide: a single analytically continued Maxwell resonance
retains a nonzero electric singular coefficient at every member of a
finite family of newly inserted conical points.  The contrast of the
insertions may depend on that finite family; the bounds
$1\leq\varepsilon_r\leq\varepsilon_+$ do not change.  A resonance here
has a complex frequency and satisfies the analytically continued
outgoing condition.  It is not a real-frequency driven solution.

\subsection{The angular branch through a homogeneous medium}

Fix an axis $e\in S^2$, a number $t\in(1/\sqrt3,1)$, and put
\begin{equation}
 a(\widehat x)=\mathbf1_{\{|\widehat x\cdot e|>t\}},\qquad
 q_\eta(\widehat x)=1+\eta a(\widehat x).
 \label{hs:modal-angular-material}
\end{equation}
The physical permittivity near a selected point is
$\varepsilon_b q_\eta$, where $1<\varepsilon_b<\varepsilon_+$ is the
permittivity of its initially homogeneous host.  The small parameter
$\eta$ is real and positive for the physical construction.  Its
complex extension below is only an analytic perturbation parameter.

\begin{lemma}[PROVED: an analytic angular eigenbranch]
\label{hs:modal-angular-branch}
For sufficiently small complex $\eta$ there are holomorphic functions
$\nu_\eta$, $\lambda_\eta$, and
$\phi_\eta\in H^1(S^2)$ satisfying
\begin{equation}
 -\operatorname{div}_{S^2}(q_\eta\nabla_{S^2}\phi_\eta)
       =\nu_\eta q_\eta\phi_\eta,\qquad
 \nu_\eta=\lambda_\eta(\lambda_\eta+1),\qquad
 \int_{S^2}q_\eta\phi_\eta^2\,dS=1.
 \label{hs:modal-angular-branch-equation}
\end{equation}
The functions $\phi_\eta$ are axisymmetric about $e$ and odd under
$\widehat x\cdot e\mapsto-\widehat x\cdot e$.  In particular their
weighted means vanish.  With the continuous choice of sign,
\begin{equation}
 \phi_0=\sqrt{\frac{3}{4\pi}}\,\widehat x\cdot e,
 \quad \nu_0=2,\quad \lambda_0=1,
 \quad \nu'_0=-3t(1-t^2),\quad
 \lambda'_0=-t(1-t^2).
 \label{hs:modal-angular-derivative}
\end{equation}
For sufficiently small real $\eta>0$, all these quantities are real
and $0<\lambda_\eta<1$.
\end{lemma}

\begin{proof}
Work in the closed subspace of axisymmetric odd functions in
$H^1(S^2)$.  At $\eta=0$, the eigenvalue two is simple in this
subspace: the spherical harmonics in it have degrees
$l=1,3,5,\ldots$, each with azimuthal index zero.  The stiffness and
mass forms have a common domain and depend affinely on $\eta$.
On this odd subspace the stiffness form is coercive at zero, and
remains invertible for small complex $\eta$ by a Neumann series.
Composing its inverse with the mass form gives a holomorphic compact
operator on $H^1$, compactness following from
$H^1(S^2)\hookrightarrow L^2(S^2)$.
A small contour about its simple eigenvalue $1/2$ defines a
holomorphic rank-one spectral projection.  Apply that projection to
$\phi_0$ and divide by the analytic square root of its nonzero
bilinear mass.  This constructs the displayed normalized eigenpair;
the analytic square root of $1+4\nu_\eta$ selects
$\lambda_\eta=(-1+\sqrt{1+4\nu_\eta})/2$ near one.
Because $q_\eta$ is axisymmetric and even, the restricted weak
eigenproblem is also the full-sphere weak eigenproblem.  Orthogonal
angular symmetry sectors make the other test functions vanish.

Differentiating the normalized eigenproblem at zero gives
\begin{equation}
 \nu'_0=\int_{S^2}a
       \bigl(|\nabla_{S^2}\phi_0|^2-2\phi_0^2\bigr)\,dS
       =3\int_t^1(1-3s^2)\,ds=-3t(1-t^2).
\end{equation}
Differentiation of $\nu=\lambda(\lambda+1)$ gives
$\lambda'_0=\nu'_0/3$.  Reality and the strict inequality follow
by conjugation uniqueness and this negative derivative.
\end{proof}

\subsection{A coefficient formula regular at the exponent one}

The apparent difficulty at $\eta=0$ is the integral
$\int_0^r s^{-\lambda_\eta}\,ds$ in the earlier coefficient
formula.  Its singular denominator is removable after one subtracts
the value of the regular curl lifting at the vertex.  The following
argument carries out that subtraction before taking the limit.

\begin{theorem}[PROVED: analytic coefficient continuation from a smooth field]
\label{hs:modal-coefficient-continuation}
Let $E_\eta$ be holomorphic with values in $L^2(B_{3R})^3$ for
$\eta$ in a complex neighborhood of zero.  Suppose that, in
$B_{3R}$, it belongs to $H_{\mathrm{loc}}(\operatorname{curl})$
and is a weak Maxwell electric field satisfying
\begin{equation}
 \operatorname{curl}\operatorname{curl}E_\eta
   =k_\eta^2\varepsilon_b q_\eta E_\eta,
 \qquad k_\eta\ne0,
 \label{hs:modal-local-maxwell}
\end{equation}
where $k_\eta$ is holomorphic and $q_\eta$ is the radially constant
coefficient in \eqref{hs:modal-angular-material}.  Then its intrinsic
conical coefficient for small real $\eta>0$ extends to a holomorphic
scalar $c_\eta$ near zero, with the exact value
\begin{equation}
 c_0=\sqrt{\frac{4\pi}{3}}\,E_0(0)\cdot e.
 \label{hs:modal-coefficient-limit}
\end{equation}
In particular, if $E_0(0)\cdot e\ne0$, then $c_\eta\ne0$ for all
sufficiently small real positive $\eta$, and the continued electric
field is unbounded at ordinary points approaching the vertex.
\end{theorem}

\begin{proof}
Fix $0<\gamma<1/2$ and decrease the parameter neighborhood so that
$\operatorname{Re}\lambda_\eta<1+\gamma/2$ and
$\operatorname{Re}\lambda_\eta>1/2$.  The local regular
decomposition of Lemma~\ref{hs:three-decomposition} also applies at
complex nonzero $k_\eta$: its cutoff estimates use only bounded
$|k_\eta|$, bounded $q_\eta$, and a positive lower bound for
$|k_\eta|$.  More explicitly, cutoff testing of
\eqref{hs:modal-local-maxwell} controls $\operatorname{curl}E_\eta$
in interior $L^2$.  Its divergence is zero and its curl is
$k_\eta^2\varepsilon_bq_\eta E_\eta$, so the constant-coefficient
curl/div identity controls it in a smaller interior $H^1$ space.
Apply one fixed regularized Poincare lifting from $H^1$ to $H^2$,
and fix the mean of the remaining scalar potential.  The result is
\begin{equation}
 E_\eta=\nabla v_\eta+W_\eta,
 \qquad v_\eta\in H^1(B_R),\quad W_\eta\in H^2(B_R)^3,
 \label{hs:modal-decomposition}
\end{equation}
holomorphically in those spaces, using slightly larger balls for
the estimates if necessary.  Holomorphy of the interior derivatives
follows from their distributional holomorphy and locally uniform
bounds; equivalently, Cauchy's integral formula in distributions
and the same local estimates gives the formula in the stronger
spaces.  Thus $W_\eta$ is holomorphic in $C^{0,\gamma}$ on the
final smaller closed ball.  From now on $R$ denotes that final
radius.

Set $w_\eta=W_\eta(0)$ and
$\widetilde W_\eta(x)=W_\eta(x)-w_\eta$.  The uniform estimate
$|\widetilde W_\eta(r\widehat x)|\leq C r^\gamma$ holds on a
smaller closed parameter disk.  Define
\begin{align}
 U_\eta(r)&=\int_{S^2}q_\eta v_\eta(r\widehat x)\phi_\eta\,dS,
 & A_{0,\eta}&=\int_{S^2}q_\eta(w_\eta\cdot\widehat x)
                                      \phi_\eta\,dS,\\
 \widetilde A_\eta(r)&=\int_{S^2}q_\eta
          (\widetilde W_\eta\cdot\widehat x)\phi_\eta\,dS,
 &\widetilde B_\eta(r)&=\int_{S^2}q_\eta
           \widetilde W_{\eta,t}\cdot\nabla_{S^2}\phi_\eta\,dS.
 \label{hs:modal-subtracted-projections}
\end{align}
All angular pairings here are bilinear, without complex conjugation,
as appropriate for holomorphic dependence.  The last two quantities
are $O(r^\gamma)$ uniformly.  For the constant vector $w_\eta$,
the eigenproblem gives the exact identity
\begin{equation}
 \int_{S^2}q_\eta w_{\eta,t}\cdot\nabla_{S^2}\phi_\eta\,dS
       =\lambda_\eta(\lambda_\eta+1)A_{0,\eta}.
 \label{hs:modal-constant-identity}
\end{equation}
Indeed $w_{\eta,t}=\nabla_{S^2}(w_\eta\cdot\widehat x)$.

Taking the divergence of \eqref{hs:modal-local-maxwell} gives
$\operatorname{div}(q_\eta E_\eta)=0$.  Its angular projection is
the same radial equation as \eqref{hs:three-radial}.  A particular
solution valid on the entire small complex parameter disk is
\begin{align}
 p_\eta(r)&=-rA_{0,\eta}+\widetilde p_\eta(r),\notag\\
 \widetilde p_\eta(r)&=\frac{1}{2\lambda_\eta+1}
 \left[r^{\lambda_\eta}\int_0^r s^{-\lambda_\eta}
       \{\widetilde B_\eta(s)
                 -(\lambda_\eta+1)\widetilde A_\eta(s)\}\,ds
 \right.\notag\\[-2pt]
 &\hspace{22mm}\left.
 -r^{-\lambda_\eta-1}\int_0^r s^{\lambda_\eta+1}
       \{\widetilde B_\eta(s)
                      +\lambda_\eta\widetilde A_\eta(s)\}\,ds
 \right].
 \label{hs:modal-regular-particular}
\end{align}
We use $r^{\lambda}=\exp(\lambda\log r)$ for $r>0$.
The integrals converge uniformly on compact parameter disks, because
their respective small-$s$ powers have real parts greater than $-1$.
They are holomorphic in $\eta$ by dominated integration, also after
parameter differentiation on a smaller disk.  Direct estimates give
\begin{equation}
 |\widetilde p_\eta(r)|\leq C r^{1+\gamma},\qquad
 |\widetilde p_\eta'(r)|\leq C r^\gamma.
 \label{hs:modal-regular-particular-bound}
\end{equation}
The first term $-rA_{0,\eta}$ is an exact particular solution for
the constant-vector forcing by \eqref{hs:modal-constant-identity}.
For $\lambda_\eta\ne1$ it is also exactly the result of performing
the two constant-vector integrals in \eqref{hs:three-exact}; thus
no additional homogeneous term has been inserted by this
regularization.

Variation of constants now gives
\begin{equation}
 U_\eta(r)=c_\eta r^{\lambda_\eta}
          +d_\eta r^{-\lambda_\eta-1}+p_\eta(r).
\end{equation}
Finite $H^1$ energy of $v_\eta$ excludes $d_\eta$: the radial
derivative of the second term would have squared integral with
weight $r^2$ proportional to
$\int_0^R r^{-2\operatorname{Re}\lambda_\eta-2}\,dr$.
The remaining coefficient has the annular formula
\begin{equation}
 c_\eta=\frac{2}{R}\int_{R/2}^R
          r^{-\lambda_\eta}\{U_\eta(r)-p_\eta(r)\}\,dr.
 \label{hs:modal-holomorphic-coefficient}
\end{equation}
The projection of $v_\eta$ on this annulus is a bounded linear
operation in $L^2$, and all other factors were shown holomorphic.
This proves holomorphy of $c_\eta$.  For positive real $\eta$,
the exponent is below one, so this is exactly the intrinsic
coefficient of Lemma~\ref{hs:three-coefficient}.

At zero the medium is homogeneous, and the Maxwell field $E_0$ is
smooth near the origin by the interior Helmholtz equation.  As
$\nabla v_0=E_0-W_0\in C^{0,\gamma}$,
\begin{equation}
 v_0(r\widehat x)=v_0(0)
        +r\{E_0(0)-w_0\}\cdot\widehat x+O(r^{1+\gamma}).
\end{equation}
The zero angular mean removes $v_0(0)$.  Adding $rA_{0,0}$ to its
projection and using the bound on $\widetilde p_0$ gives
\begin{equation}
 c_0=\int_{S^2}(E_0(0)\cdot\widehat x)\phi_0\,dS
     =\sqrt{\frac{4\pi}{3}}\,E_0(0)\cdot e.
\end{equation}
The nonvanishing and unboundedness conclusions follow from
continuity and Lemma~\ref{hs:three-coefficient}.  In particular,
closeness of fields in $L^2$ alone is not being used to infer
nonvanishing of a singular coefficient; its explicit continuous
extension has been established.
\end{proof}

\begin{corollary}[PROVED: all finitely many coefficients of one continued mode]
\label{hs:modal-all-coefficients}
Suppose a compact dielectric has a Maxwell resonance $k_0\ne0$
with eigenfield $E_0$.  Let $p_1,\ldots,p_N$ have disjoint
neighborhoods of constant permittivity $\varepsilon_b$, and choose
axes $e_j$ such that $E_0(p_j)\cdot e_j\ne0$.
In each neighborhood insert the angular material
$\varepsilon_b q_\eta$ of
\eqref{hs:modal-angular-material}, with that neighborhood's axis,
and close the insertions outside its local conical ball.
If the same resonance has a nonzero local-$L^2$ holomorphic
continuation $E_\eta,k_\eta$ for this simultaneous material
perturbation, with $E_\eta\vert_{\eta=0}=E_0$, then there is
$\eta_N>0$ such that every $0<\eta<\eta_N$ has
\begin{equation}
 c_j(E_\eta)\ne0\quad\hbox{for all }1\leq j\leq N.
 \label{hs:modal-all-coefficients-equation}
\end{equation}
These are coefficients of one eigenfield at one resonance, not of
a sum of independent modes or separately selected illuminations.
The physical material bound is respected by also requiring
$\varepsilon_b(1+\eta)\leq\varepsilon_+$.
\end{corollary}

\begin{proof}
Translate and rotate Theorem~\ref{hs:modal-coefficient-continuation}
at each $p_j$ and take the minimum of the finitely many positive
nonvanishing radii in parameter space.  Scaling a chosen eigenfield
by any holomorphic nonzero normalization factor multiplies all its
coefficients by that same factor and does not affect the result.
\end{proof}

\subsection{The exact remaining global hypotheses}

\begin{proposition}[PROVED: transfer from a simple resonance to rounded modal peaks]
\label{hs:modal-transfer}
Assume a finite homogeneous dielectric sphere with permittivity
$\varepsilon_b\in(1,\varepsilon_+)$ has a simple resonance in the
axisymmetric TM sector, whose axial electric component is nonzero
on an open interval of its symmetry axis.  Assume furthermore:
\begin{enumerate}
\item This sector's eigenpair has a holomorphic continuation under
the finite axisymmetric cone insertions in
Corollary~\ref{hs:modal-all-coefficients}.
\item An explicitly specified genuine QNM normalization is finite,
nonzero and continuous along that branch and under the rounded
approximations below.
\item Rounding the conical vertices while retaining the same
material values gives a nearby resonance and normalized eigenfield
converging locally in $L^2$, and uniformly one-sided on fixed
smooth-interface collars away from the vertices.
\end{enumerate}
Then, for every finite $N$ and every $M,\ell>0$, one normalized
resonance of a finite smooth dielectric has $N$ disjoint
three-dimensional peak balls of diameter less than $\ell$, each
with one-sided maximum larger than $M$ and a collar maximum smaller
than half its peak.  Their centers may be placed on a line in a
prescribed open plane patch, and their projections may be contained
in that patch, provided that the chosen nonzero axial interval is
positioned to meet the patch.  If the symmetry axis lies in that
plane and the rounding remains axisymmetric, a maximizing point
in each ball can itself be selected in the plane patch.  The
surrounding resonator need not project into the patch.  A maximum
off the symmetry axis may belong to a rotational ring inside its
small ball; no nondegeneracy or uniqueness of a maximum is asserted.
\end{proposition}

\begin{proof}
Choose $N$ distinct axial centers in a sufficiently short nonzero
subinterval and choose disjoint balls and insertions as small as
necessary.  Axisymmetric insertions preserve the indicated sector.
Choose one small positive contrast from
Corollary~\ref{hs:modal-all-coefficients}.  Its one normalized sharp
resonance is unbounded at every tip and bounded on each chosen
fixed collar.  Select an ordinary high-field point inside each ball
whose value exceeds twice the collar bound and the target $M$, with
strict margins.  The hypothesized normalized rounding convergence
preserves all these finitely many strict inequalities at one
sufficiently small rounding scale.  Smooth-interface regularity
makes the one-sided maxima on the closed balls finite and attained;
each peak lies inside its ball because of the collar inequality.
This is the same final compactness step as the driven construction,
but it is applied to a homogeneous resonance branch with its stated
QNM normalization.

For the last placement assertion, an axisymmetric vector field
obeys $E(\mathcal R x)=\mathcal R E(x)$ under rotations about its
axis, and therefore $|E|$ is invariant.  The same is true for its
one-sided material traces.  A ball centered on that axis is
rotation invariant.  Rotate any maximizing point into the specified
plane through the axis; it stays in the same ball and retains its
maximal value.  The projection constraint puts this representative
in the prescribed patch.
\end{proof}

This proposition is a proved implication with explicit global
hypotheses.  Lemmas~\ref{ms:sphere-simple}, \ref{ms:continuation},
\ref{ms:normalization}, and~\ref{ms:rounding} supply those hypotheses
for the finite-device construction in Theorem~\ref{ms:modal-packing}.
Theorem~\ref{hs:modal-coefficient-continuation} and
Corollary~\ref{hs:modal-all-coefficients} establish its local
excitation bridge.  No resonance is inferred from a driven field,
and no radiating mode is normalized by a divergent all-space positive
energy integral.  A real dilation fixes the real part of the
nonzero resonance frequency; its negative imaginary part remains
part of the QNM statement.
